% =========================================================
%  ppt/05-sprite-group-collision.tex
%  Materi 5 — Sprite, Group, dan Collision Detection
%  Bab VII buku ajar
% =========================================================
\documentclass[aspectratio=169,10pt]{beamer}
% =========================================================
%  ppt/config/packages.tex
%  Paket bersama untuk seluruh presentasi kuliah
%  Pengantar Pemrograman Game
% =========================================================

\usepackage[T1]{fontenc}
\usepackage[utf8]{inputenc}
\usepackage[indonesian]{babel}
\usepackage{lmodern}
\usepackage{textcomp}
\usepackage{amsmath}
\usepackage{amssymb}
\usepackage{graphicx}
\usepackage{booktabs}
\usepackage{tabularx}
\usepackage{array}
\usepackage{multicol}
\usepackage{listings}
\usepackage{upquote}
\usepackage{adjustbox}
\usepackage{tcolorbox}
\usepackage{tikz}
\usetikzlibrary{
  arrows.meta,
  positioning,
  shapes.geometric,
  shapes.misc,
  calc,
  fit,
  backgrounds,
  matrix,
  chains,
  decorations.pathreplacing,
  patterns
}

\tcbuselibrary{skins}

% Jangan memuat geometry: ukuran halaman diatur oleh beamer.

% =========================================================
%  ppt/config/theme.tex
%  Satu tema Beamer untuk keempat belas presentasi
%  Pengantar Pemrograman Game
% =========================================================

% --------------------------------------------------------- warna
\definecolor{ubNavy}{HTML}{1F3864}
\definecolor{ubAccent}{HTML}{2E75B6}
\definecolor{ubHead}{HTML}{D9E2F3}
\definecolor{ubSoft}{HTML}{F2F6FC}
\definecolor{ubLine}{HTML}{BFBFBF}
\definecolor{ubText}{HTML}{1F2937}
\definecolor{ubMuted}{HTML}{5A6678}
\definecolor{ubGreen}{HTML}{1E7B45}
\definecolor{ubGreenBg}{HTML}{F1F8F2}
\definecolor{ubRed}{HTML}{B3261E}
\definecolor{ubRedBg}{HTML}{FDF2F1}
\definecolor{ubWarn}{HTML}{A35A05}
\definecolor{ubWarnBg}{HTML}{FFF8EC}
\definecolor{ubGray}{HTML}{F6F6F6}
\definecolor{ubCodeBg}{HTML}{10233F}
\definecolor{ubCodeText}{HTML}{E8EEF7}
\definecolor{ubCodeKey}{HTML}{8FBCF0}
\definecolor{ubCodeCom}{HTML}{8FB37A}
\definecolor{ubCodeStr}{HTML}{D9A46B}
\definecolor{ubCodeNum}{HTML}{7E99BE}

% Nama warna yang dipakai gambar TikZ buku ajar.
\definecolor{ubPrimary}{HTML}{1F3864}
\definecolor{ubAlt}{HTML}{F2F6FC}

% --------------------------------------------------------- dasar
\usetheme{default}
\usefonttheme{professionalfonts}
\setbeamertemplate{navigation symbols}{}
\setbeamerfont{frametitle}{size=\Large,series=\bfseries}
\setbeamerfont{framesubtitle}{size=\small}
\setbeamerfont{footline}{size=\tiny}
\setbeamerfont{block title}{size=\small,series=\bfseries}
\setbeamerfont{block body}{size=\small}
\setbeamerfont{itemize/enumerate body}{size=\small}
\setbeamerfont{itemize/enumerate subbody}{size=\footnotesize}

\setbeamersize{text margin left=0.60cm,text margin right=0.60cm}
\uchyph=0

\setbeamercolor{normal text}{fg=ubText,bg=white}
\setbeamercolor{structure}{fg=ubAccent}
\setbeamercolor{frametitle}{fg=ubNavy,bg=white}
\setbeamercolor{title}{fg=ubNavy}
\setbeamercolor{footline}{fg=ubMuted,bg=white}
\setbeamercolor{itemize item}{fg=ubAccent}
\setbeamercolor{itemize subitem}{fg=ubAccent}
\setbeamercolor{block title}{fg=white,bg=ubNavy}
\setbeamercolor{block body}{fg=ubText,bg=ubSoft}
\setbeamercolor{block title alerted}{fg=white,bg=ubWarn}
\setbeamercolor{block body alerted}{fg=ubText,bg=ubWarnBg}
\setbeamercolor{block title example}{fg=white,bg=ubGreen}
\setbeamercolor{block body example}{fg=ubText,bg=ubGreenBg}
\setbeamercolor{caption name}{fg=ubMuted}

\setbeamertemplate{blocks}[rounded][shadow=false]
\setbeamertemplate{itemize item}{\raisebox{0.10em}{\color{ubAccent}\scriptsize$\blacksquare$}}
\setbeamertemplate{itemize subitem}{\raisebox{0.06em}{\color{ubAccent}\tiny$\blacktriangleright$}}
\setbeamertemplate{enumerate item}{%
  \makebox[1.3em][l]{\color{ubAccent}\insertenumlabel.}}

\addtobeamertemplate{itemize/enumerate body begin}{%
  \setlength{\itemsep}{0.42em}\setlength{\parskip}{0pt}}{}

% --------------------------------------------------------- judul frame
\setbeamertemplate{frametitle}[default][left]

% --------------------------------------------------------- kaki
\newcommand{\materino}{00}
\newcommand{\materifoot}{Pengantar Pemrograman Game \textperiodcentered\ Materi \materino}

\setbeamertemplate{footline}{%
  \begin{beamercolorbox}[wd=\paperwidth,ht=0.015cm,dp=0cm]{footline}%
    \color{ubHead}\rule{\paperwidth}{0.4pt}%
  \end{beamercolorbox}%
  \begin{beamercolorbox}[wd=\paperwidth,ht=0.40cm,dp=0.14cm,%
                         leftskip=0.60cm,rightskip=0.60cm]{footline}%
    \usebeamerfont{footline}\color{ubMuted}\materifoot\hfill\insertframenumber%
  \end{beamercolorbox}%
}

% --------------------------------------------------------- makro buku ajar
\providecommand{\code}[1]{\texttt{\detokenize{#1}}}
\providecommand{\istilah}[1]{\emph{#1}}
\providecommand{\CPMK}[1]{CPMK\,#1}
\providecommand{\CPL}[1]{CPL\,#1}
\providecommand{\SUBCPMK}[1]{Sub-CPMK\,#1}

% --------------------------------------------------------- gaya kode
\lstdefinelanguage{CSharp}{
  morekeywords={abstract,as,base,bool,break,byte,case,catch,char,class,const,
    continue,decimal,default,delegate,do,double,else,enum,event,explicit,
    extern,false,finally,fixed,float,for,foreach,goto,if,implicit,in,int,
    interface,internal,is,lock,long,namespace,new,null,object,operator,out,
    override,params,private,protected,public,readonly,ref,return,sbyte,
    sealed,short,sizeof,stackalloc,static,string,struct,switch,this,throw,
    true,try,typeof,uint,ulong,unchecked,unsafe,ushort,using,virtual,void,
    volatile,while,var,get,set},
  sensitive=true,
  morecomment=[l]{//},
  morecomment=[s]{/*}{*/},
  morestring=[b]",
}

\lstdefinestyle{kode}{
  basicstyle=\ttfamily\footnotesize\color{ubCodeText},
  keywordstyle=\color{ubCodeKey}\bfseries,
  commentstyle=\color{ubCodeCom}\itshape,
  stringstyle=\color{ubCodeStr},
  numberstyle=\tiny\color{ubCodeNum},
  numbers=left,
  numbersep=7pt,
  backgroundcolor=\color{ubCodeBg},
  frame=none,
  showstringspaces=false,
  breaklines=false,
  tabsize=4,
  columns=fullflexible,
  keepspaces=true,
  xleftmargin=12pt,
  xrightmargin=2pt,
  aboveskip=0pt,
  belowskip=0pt,
  upquote=true,
}
\lstset{style=kode}

\newcommand{\kodejelas}[1]{%
  \vspace{0.16cm}%
  {\footnotesize\color{ubMuted}#1\par}%
}
\newcommand{\sumber}[1]{%
  \vspace{0.08cm}%
  {\scriptsize\color{ubMuted}\itshape #1\par}%
}
\newcommand{\awalan}[1]{%
  {\small\color{ubMuted}#1\par}\vspace{0.18cm}%
}

% --------------------------------------------------------- kotak
\newtcolorbox{kotak}[1][]{enhanced, colback=ubSoft, colframe=ubAccent,
  boxrule=0pt, leftrule=2.2pt, arc=1pt, left=6pt, right=6pt, top=4pt,
  bottom=4pt, fontupper=\footnotesize, #1}
\newtcolorbox{peringatan}[1][]{enhanced, colback=ubWarnBg, colframe=ubWarn,
  boxrule=0pt, leftrule=2.2pt, arc=1pt, left=6pt, right=6pt, top=4pt,
  bottom=4pt, fontupper=\footnotesize, #1}
\newtcolorbox{kiat}[1][]{enhanced, colback=ubGreenBg, colframe=ubGreen,
  boxrule=0pt, leftrule=2.2pt, arc=1pt, left=6pt, right=6pt, top=4pt,
  bottom=4pt, fontupper=\footnotesize, #1}

% --------------------------------------------------------- sampul
% \deckcover{nomor}{minggu}{judul}{sub-cpmk}{teknologi}{sumber}{blok}
\newcommand{\deckcover}[7]{%
\begin{frame}[plain]
\begin{tikzpicture}[remember picture, overlay]
  \fill[ubNavy] (current page.south west) rectangle ([xshift=6.00cm]current page.north west);
  \fill[ubAccent] ([xshift=6.00cm]current page.south west) rectangle ([xshift=6.08cm]current page.north west);

  \node[anchor=north west, text width=4.72cm, align=left, inner sep=0pt]
    at ($(current page.north west)+(0.66cm,-0.60cm)$) {%
      {\color{white}\bfseries\fontsize{10.2}{12.2}\selectfont
       Universitas Bale Bandung}\\[4pt]
      {\color{ubHead}\fontsize{7.6}{9.4}\selectfont
       Fakultas Teknologi Informasi}\\[1pt]
      {\color{ubHead}\fontsize{7.6}{9.4}\selectfont
       Program Studi Teknik Informatika}\\[7pt]
      {\color{ubAccent}\rule{4.0cm}{0.5pt}}%
    };

  \node[anchor=north west, text width=4.72cm, align=left, inner sep=0pt]
    at ($(current page.north west)+(0.66cm,-2.72cm)$) {%
      {\color{ubAccent}\fontsize{6.4}{7.8}\selectfont MATA KULIAH}\\[3pt]
      {\color{white}\bfseries\fontsize{9.8}{12.0}\selectfont
       Pengantar Pemrograman Game}\\[3pt]
      {\color{ubHead}\fontsize{7.6}{9.4}\selectfont
       MWCA009 \textperiodcentered\ 3 SKS}\\[8pt]
      {\color{ubHead}\fontsize{7.6}{9.4}\selectfont
       Semester Ganjil 2026/2027}%
    };

  \node[anchor=south west, text width=4.72cm, align=left, inner sep=0pt]
    at ($(current page.south west)+(0.66cm,0.52cm)$) {%
      {\color{ubAccent}\fontsize{6.4}{7.8}\selectfont DOSEN PENGAMPU}\\[3pt]
      {\color{white}\fontsize{8.4}{10.4}\selectfont
       Cecep Suwanda, S.Si., M.Kom.}%
    };

  \node[anchor=north west, text width=8.95cm, align=left, inner sep=0pt]
    at ($(current page.north east)+(-9.40cm,-0.92cm)$) {%
      {\color{ubAccent}\bfseries\fontsize{8.2}{10.2}\selectfont
       MATERI #1\ \ \textperiodcentered\ \ MINGGU #2}\\[8pt]
      {\color{ubNavy}\bfseries\fontsize{19.4}{23.4}\selectfont
       \hyphenpenalty=10000 \exhyphenpenalty=10000 \raggedright #3}%
    };

  \node[anchor=north west, text width=8.95cm, align=left, inner sep=0pt]
    at ($(current page.north east)+(-9.40cm,-4.30cm)$) {%
      \renewcommand{\arraystretch}{1.30}%
      \begin{tabular}{@{}>{\color{ubMuted}\fontsize{7.3}{9.0}\selectfont}p{1.85cm}
                        >{\color{ubText}\fontsize{8.1}{10.0}\selectfont}p{6.7cm}@{}}
        Sub-CPMK      & #4 \\
        Teknologi     & #5 \\
        Sumber materi & #6 \\
        Blok          & #7 \\
      \end{tabular}%
    };
\end{tikzpicture}
\end{frame}%
}

% --------------------------------------------------------- pembatas bagian
\newcommand{\sectionframe}[3]{%
\begin{frame}[plain]
\begin{tikzpicture}[remember picture, overlay]
  \fill[ubNavy] (current page.south west) rectangle (current page.north east);
  \fill[ubAccent] (current page.north west) rectangle ([xshift=0.10cm]current page.south west);
  \node[anchor=west, inner sep=0pt]
    at ($(current page.west)+(1.00cm,1.45cm)$)
    {\color{ubNavy!52!white}\bfseries\fontsize{58}{58}\selectfont #1};
  \node[anchor=west, inner sep=0pt, text width=11.6cm, align=left]
    at ($(current page.west)+(1.05cm,-0.55cm)$) {%
      {\color{ubAccent}\bfseries\fontsize{8.0}{10.0}\selectfont BAGIAN}\\[3pt]
      {\color{white}\bfseries\fontsize{20.5}{24.5}\selectfont #2}\\[5pt]
      {\color{ubHead}\small #3}%
    };
  \node[anchor=south east, inner sep=0pt]
    at ($(current page.south east)+(-0.60cm,0.32cm)$)
    {\color{ubHead}\fontsize{7.4}{9}\selectfont #1 / 6
     \ \ \textperiodcentered\ \ \insertframenumber};
\end{tikzpicture}
\end{frame}%
}

% --------------------------------------------------------- peta pertemuan
\newcommand{\blokpeta}[6]{%
  \ifnum#2=#3
    \node[rounded corners=3pt, fill=ubNavy, draw=ubNavy, line width=0.5pt,
          text width=3.02cm, minimum height=1.62cm, align=left, inner sep=5pt]
      at (#1,0) {%
        {\color{ubAccent}\bfseries\fontsize{6.8}{8}\selectfont #4}\\[1pt]
        {\color{white}\bfseries\fontsize{7.6}{9}\selectfont #5}\\[2pt]
        {\color{ubHead}\fontsize{6.4}{7.6}\selectfont #6}};
  \else
    \node[rounded corners=3pt, fill=ubSoft, draw=ubHead, line width=0.5pt,
          text width=3.02cm, minimum height=1.62cm, align=left, inner sep=5pt]
      at (#1,0) {%
        {\color{ubAccent}\bfseries\fontsize{6.8}{8}\selectfont #4}\\[1pt]
        {\color{ubNavy}\bfseries\fontsize{7.6}{9}\selectfont #5}\\[2pt]
        {\color{ubMuted}\fontsize{6.4}{7.6}\selectfont #6}};
  \fi
}

% \petapertemuan{blok aktif}{sebelumnya}{sesudahnya}{keluaran}
% \petapertemuan{blok aktif}{sebelumnya}{pertemuan ini}{sesudahnya}{keluaran}
\newcommand{\petapertemuan}[5]{%
\begin{frame}{Peta Pertemuan}
{\scriptsize\color{ubMuted}Kedudukan materi dalam empat blok perkuliahan}\par
\vspace{0.10cm}
\centering
\begin{tikzpicture}
  \blokpeta{-5.58}{1}{#1}{BLOK I}{Materi 1--5}{Game loop, vektor,\\fisika, collision}
  \blokpeta{-1.86}{2}{#1}{BLOK II}{Materi 6--9}{AI NPC: chase, FSM,\\pathfinding A*}
  \blokpeta{1.86}{3}{#1}{BLOK III}{Materi 10--11}{Unity: scene,\\level, physics}
  \blokpeta{5.58}{4}{#1}{BLOK IV}{Materi 12--14}{Aset visual, audio,\\UI, playtest}
\end{tikzpicture}
\par\vspace{0.30cm}
\begin{columns}[T,onlytextwidth]
  \begin{column}{0.315\textwidth}
    \begin{beamercolorbox}[wd=\linewidth,sep=4pt,rounded=true]{block body}
      {\color{ubMuted}\fontsize{6.6}{8}\selectfont SEBELUMNYA}\\[2pt]
      {\footnotesize #2}
    \end{beamercolorbox}
  \end{column}
  \begin{column}{0.315\textwidth}
    \begin{beamercolorbox}[wd=\linewidth,sep=4pt,rounded=true]{block title}
      {\color{ubHead}\fontsize{6.6}{8}\selectfont PERTEMUAN INI}\\[2pt]
      {\footnotesize #3}
    \end{beamercolorbox}
  \end{column}
  \begin{column}{0.315\textwidth}
    \begin{beamercolorbox}[wd=\linewidth,sep=4pt,rounded=true]{block body}
      {\color{ubMuted}\fontsize{6.6}{8}\selectfont SESUDAHNYA}\\[2pt]
      {\footnotesize #4}
    \end{beamercolorbox}
  \end{column}
\end{columns}
\par\vspace{0.26cm}
{\footnotesize\color{ubText}\textbf{Keluaran pertemuan.} #5\par}
\end{frame}%
}

% --------------------------------------------------------- penutup
\newcommand{\penutup}[2]{%
\begin{frame}[plain]
\begin{tikzpicture}[remember picture, overlay]
  \fill[ubNavy] (current page.south west) rectangle (current page.north east);
  \fill[ubAccent] (current page.north west) rectangle ([xshift=0.10cm]current page.south west);

  \node[anchor=north west, text width=13.6cm, align=left, inner sep=0pt]
    at ($(current page.north west)+(1.00cm,-0.95cm)$) {%
      {\color{ubAccent}\bfseries\fontsize{8.2}{10}\selectfont ASESMEN MINGGU INI}\\[6pt]
      {\color{white}\bfseries\fontsize{24}{28}\selectfont #1}%
    };

  \node[anchor=north west, text width=13.6cm, align=left, inner sep=0pt]
    at ($(current page.north west)+(1.00cm,-4.55cm)$) {%
      {\color{ubNavy!0!ubHead}\rule{13.4cm}{0.4pt}}\\[7pt]
      {\color{ubAccent}\bfseries\fontsize{7.4}{9}\selectfont
       PENGAIT KE MATERI BERIKUTNYA}\\[4pt]
      {\color{ubHead}\fontsize{10.6}{13.2}\selectfont #2}%
    };

  \node[anchor=south east, inner sep=0pt]
    at ($(current page.south east)+(-0.60cm,0.32cm)$)
    {\color{ubHead}\fontsize{7.4}{9}\selectfont \insertframenumber};
\end{tikzpicture}
\end{frame}%
}



\renewcommand{\materino}{05}

\begin{document}

\deckcover{05}{5}
  {Sprite, Group, dan Collision Detection}
  {1.4}
  {Python / Pygame}
  {Bab VII buku ajar}
  {Blok I \textperiodcentered\ Materi 1--5}

% --------------------------------------------------------- capaian
\begin{frame}{Capaian Pertemuan}
\setbeamerfont{itemize/enumerate body}{size=\scriptsize}
\awalan{Sub-CPMK dan indikator dari RPS, serta tujuan Bab VII buku ajar tentang
\emph{collision detection}.}
\begin{columns}[T,onlytextwidth]
  \begin{column}{0.52\textwidth}
    \begin{beamercolorbox}[wd=\linewidth,sep=5pt,rounded=true]{block title}
      {\fontsize{6.8}{8}\selectfont\color{ubHead}SUB-CPMK 1.4}\\[2pt]
      {\footnotesize Mendeteksi dan menangani \emph{collision} menggunakan
       \texttt{Rect} dan/atau \emph{mask} pada \emph{sprite}.}
    \end{beamercolorbox}
    \vspace{0.16cm}

    {\fontsize{7}{8.5}\selectfont\bfseries\color{ubAccent}INDIKATOR}\par
    \vspace{0.06cm}
    \begin{itemize}
      \item Memilih metode \emph{collision} yang tepat sesuai kebutuhan.
      \item Menyelesaikan tabrakan per sumbu sehingga objek tidak menembus
            dinding.
      \item Menganalisis dan memperbaiki kasus tabrakan yang salah
            (tumpang tindih atau lolos deteksi).
    \end{itemize}
  \end{column}
  \begin{column}{0.45\textwidth}
    {\fontsize{7}{8.5}\selectfont\bfseries\color{ubAccent}TUJUAN BAB VII}\par
    \vspace{0.08cm}
    \begin{itemize}
      \item Menyusun benda sebagai kelas turunan
            \texttt{pygame.sprite.Sprite} dan mengelolanya dengan
            \texttt{Group}.
      \item Memilih di antara \texttt{colliderect}, \texttt{spritecollide},
            dan \texttt{groupcollide}.
      \item Menyelesaikan tabrakan sumbu demi sumbu sehingga benda dapat
            mendarat sekaligus terhalang.
      \item Menjelaskan kapan persegi pembatas tidak memadai dan bagaimana
            \emph{mask} dipakai.
      \item Menganalisis tiga gejala tabrakan yang salah.
      \item Memeriksa kebenaran penanganan tabrakan secara otomatis.
    \end{itemize}
  \end{column}
\end{columns}
\end{frame}

% --------------------------------------------------------- peta
\petapertemuan{1}
  {Materi 4 \textperiodcentered\ Fisika Sederhana}
  {Materi 5 \textperiodcentered\ Sprite, Group, dan Collision Detection}
  {Materi 6 \textperiodcentered\ Pengantar AI NPC dan Perilaku Chase}
  {Prototipe permainan yang pemainnya mendarat, terhalang, mengumpulkan
   kepingan, dan menghindari rintangan; penghitung tumpang tindih bernilai nol.}

% ========================================================= bagian 1
\section{Kelas Sprite dan Wadah Group}
\sectionframe{01}{Kelas Sprite dan Wadah Group}
  {Satu kelas menghasilkan banyak benda, satu wadah mengurus semuanya}

\begin{frame}{Mengapa Benda Perlu Menjadi Satu Kesatuan}
\awalan{Pada akhir Bab VI, setiap benda masih disimpan pada beberapa variabel
terpisah. Cara itu tidak lagi memadai ketika jumlah benda bertambah.}
\begin{itemize}
  \item Pemain disimpan pada empat variabel: posisi, kecepatan, \texttt{rect},
        dan penanda menyentuh lantai; bola pada tiga variabel serupa.
  \item Menambah benda ketiga berarti menambah tiga atau empat variabel lagi,
        beserta baris pembaruannya.
  \item Tiga kesulitan muncul: jumlah variabel bertambah secepat jumlah benda;
        mudah keliru memasangkan variabel; dan perintah untuk seluruh benda
        harus ditulis berulang.
\end{itemize}
\vspace{0.16cm}
\begin{kotak}
Pemrograman berorientasi objek mengumpulkan data dan perilaku sebuah benda
menjadi satu objek, dengan kelas sebagai resep pembuatnya. Fungsi
\texttt{\_\_init\_\_} dijalankan sekali saat objek dibuat, dan \texttt{self}
menunjuk objek yang sedang dikerjakan.
\end{kotak}
\end{frame}

\begin{frame}[fragile]{Pola Kelas \emph{Sprite} pada Pygame}
\awalan{Pygame menyediakan kelas dasar \texttt{pygame.sprite.Sprite}. Tugas Anda
menurunkannya dan mengisi gambar serta persegi pembatasnya.}
\lstinputlisting[language=Python,firstline=81,lastline=88]
  {../buku_ajar/code/python/ch07_tabrakan.py}
\kodejelas{\texttt{super().\_\_init\_\_()} menjalankan penyiapan kelas dasar dan
wajib ada. Nama \texttt{self.image} dan \texttt{self.rect} bukan pilihan Anda,
melainkan kesepakatan Pygame: perkakas penggambaran dan pemeriksaan tabrakan
mencari kedua nama itu.}
\vspace{0.15cm}
\begin{peringatan}
Bila keduanya diberi nama lain, program tetap berjalan tetapi benda tidak
tergambar atau tidak pernah dianggap bertabrakan. Gejalanya membingungkan
karena tidak disertai pesan kesalahan.
\end{peringatan}
\end{frame}

\begin{frame}{Wadah \emph{Group}}
\begin{center}
\adjustbox{max width=0.74\linewidth,max height=0.43\textheight}{%
  % chapters/07/figures/sprite-group.tikz
% Hubungan kelas Sprite dan wadah Group: satu kelas menghasilkan banyak benda,
% dan wadah mengerjakan perintah untuk seluruh bendanya sekaligus.
\begin{tikzpicture}[
  kotak/.style={
    rectangle, rounded corners=2pt, draw=ubPrimary, thick, fill=ubAlt,
    align=center, inner sep=4pt, font=\footnotesize
  },
  wadah/.style={
    rectangle, rounded corners=2pt, draw=ubAccent, thick, fill=white,
    align=center, inner sep=6pt, font=\footnotesize
  },
  benda/.style={circle, draw=ubAccent, thick, fill=ubAccent!25, inner sep=3pt},
  panah/.style={-{Stealth[length=2.4mm]}, thick, draw=ubAccent},
  label/.style={font=\scriptsize, align=center}
]

% Kelas sebagai resep, di sebelah kiri.
\node[kotak, text width=34mm] (kelas) at (0,0)
  {\textbf{kelas \code{Item}}\\[2pt]\scriptsize resep benda: gambar,\\rect, dan perilakunya};

% Tiga benda hasil kelas tersebut, di dalam wadah Group di sebelah kanan.
\node[wadah, text width=48mm, minimum height=20mm] (grup) at (7.6,0) {};
\node[anchor=south, font=\scriptsize\bfseries, ubAccent] at (7.6,1.15)
  {wadah \code{Group}};

\node[benda] at (6.4,-0.1) {};
\node[benda] at (7.6,-0.1) {};
\node[benda] at (8.8,-0.1) {};

\draw[panah] (kelas.east) -- node[label, above, yshift=1pt]
  {membuat tiga benda} (grup.west);

% Keterangan di bawah diagram, pada barisnya sendiri.
\node[label, anchor=north, text width=124mm] at (3.8,-1.7)
  {\code{update()} dan \code{draw()} diberikan kepada wadah, lalu wadah
   mengerjakannya untuk setiap benda di dalamnya. Perintah tidak perlu ditulis
   satu per satu.};

\node[label, anchor=north, text width=124mm] at (3.8,-2.7)
  {Satu kelas dipakai berulang kali. Menambah benda baru cukup dengan membuat
   satu benda lagi dan memasukkannya ke dalam wadah.};
\end{tikzpicture}
}
\end{center}
\vspace{-0.06cm}
\begin{itemize}
  \item \texttt{pygame.sprite.Group} menyimpan sekumpulan benda dan
        mengerjakan \texttt{update} serta \texttt{draw} untuk seluruhnya
        sekaligus.
  \item Satu benda boleh berada pada beberapa wadah; wadah hanyalah kumpulan
        rujukan, bukan salinan.
  \item \texttt{group.add(benda)} memasukkan anggota, \texttt{sprite.kill()}
        mengeluarkannya dari seluruh wadah, dan \texttt{for benda in group:}
        menelusuri isinya seperti daftar biasa.
\end{itemize}
\end{frame}

% ========================================================= bagian 2
\section{Deteksi Tabrakan dengan Persegi Pembatas}
\sectionframe{02}{Deteksi Tabrakan dengan Persegi Pembatas}
  {Tiga perintah, tiga keperluan yang berbeda}

\begin{frame}{Tiga Perintah Pemeriksaan Tabrakan}
\awalan{Seluruh perintah berikut disediakan Pygame dan bekerja pada persegi
pembatas (\emph{AABB}) yang disimpan pada \texttt{rect}.}
\begin{itemize}
  \item \texttt{rect.colliderect(other)} memeriksa satu persegi terhadap satu
        persegi. Dua persegi yang hanya bersentuhan pada tepinya dinilai
        \textbf{tidak} beririsan.
  \item \texttt{spritecollide(benda, wadah, dokill)} memeriksa satu benda
        terhadap sekumpulan benda dan mengembalikan daftar.
  \item \texttt{groupcollide(wadah1, wadah2, kill1, kill2)} memeriksa dua
        kumpulan dan mengembalikan kamus, bukan daftar.
\end{itemize}
\vspace{0.16cm}
\begin{kotak}
Sifat ``tepi yang bersentuhan tidak beririsan'' berguna: benda yang sudah
ditempelkan pada tepi bongkahan tidak dianggap bertabrakan lagi pada bingkai
berikutnya.
\end{kotak}
\end{frame}

\begin{frame}[fragile]{Memakainya pada Program}
\lstinputlisting[language=Python,firstline=249,lastline=262]
  {../buku_ajar/code/python/ch07_tabrakan.py}
\kodejelas{Argumen ketiga \texttt{spritecollide} menentukan apakah benda yang
tersentuh langsung dikeluarkan dari wadah. Nilai benar dipakai untuk kepingan
yang diambil; nilai salah dipakai untuk rintangan yang harus tetap di
tempatnya. Jumlah kepingan dihitung dengan \texttt{len(hit\_items)}, sebab satu
bingkai dapat menyentuh lebih dari satu kepingan.}
\end{frame}

\begin{frame}{Mendeteksi Bukan Menyelesaikan}
\begin{itemize}
  \item Seluruh perintah pada bagian sebelumnya hanya menjawab pertanyaan
        ``apakah kedua benda bersentuhan?''.
  \item Pertanyaan berikutnya---``dari sisi mana tabrakan terjadi, dan apa yang
        harus dilakukan?''---belum terjawab.
  \item Jawaban itu bergantung pada aturan permainan Anda, bukan pada pustaka:
        pada satu permainan menyentuh kepingan menambah skor, pada permainan
        lain menyentuh benda yang sama mengurangi nyawa.
\end{itemize}
\vspace{0.2cm}
\begin{kotak}
Mendeteksi tabrakan dan menyelesaikannya adalah dua pekerjaan yang berbeda.
Pustaka menjawab yang pertama; aturan permainan Anda menentukan yang kedua.
Bagian berikutnya membahas pekerjaan kedua itu.
\end{kotak}
\end{frame}

% ========================================================= bagian 3
\section{Menyelesaikan Tabrakan Sumbu demi Sumbu}
\sectionframe{03}{Menyelesaikan Tabrakan Sumbu demi Sumbu}
  {Satu sumbu diselesaikan tuntas sebelum sumbu berikutnya disentuh}

\begin{frame}{Mengapa Perlu Dipisah per Sumbu}
\begin{center}
\adjustbox{max width=0.88\linewidth,max height=0.42\textheight}{%
  % chapters/07/figures/tabrakan-per-sumbu.tikz
% Penyelesaian tabrakan sumbu demi sumbu: gerak mendatar diperiksa lebih
% dahulu, baru gerak tegak.
\begin{tikzpicture}[
  bongkah/.style={rectangle, draw=ubLine, thick, fill=ubGray},
  pemain/.style={rectangle, draw=ubPrimary, thick, fill=ubPrimary!25},
  pemainlama/.style={rectangle, draw=ubPrimary, dashed, thin},
  panah/.style={-{Stealth[length=2.2mm]}, thick, draw=ubAccent},
  judul/.style={font=\footnotesize\bfseries, ubPrimary, align=center},
  label/.style={font=\scriptsize, align=left}
]

% ---------- Panel A: sumbu mendatar ----------
\begin{scope}
  \node[judul, anchor=south] at (2.4,2.15)
    {1. Sumbu mendatar diperiksa lebih dahulu};

  % Bongkahan
  \draw[bongkah] (3.2,0) rectangle (5.6,2.0);
  \node[label, anchor=north] at (4.4,-0.15) {bongkahan};

  % Posisi sebelum dan sesudah penyelesaian
  \draw[pemainlama] (2.55,0.55) rectangle (3.05,1.45);
  \draw[pemain] (2.7,0.55) rectangle (3.2,1.45);

  \draw[panah] (2.2,1.7) -- (3.0,1.7);
  \node[label, anchor=south west] at (2.25,1.75)
    {bergerak ke kanan};

  % Tanda sisi yang ditempelkan
  \draw[ubRed, very thick] (3.2,0.4) -- (3.2,1.6);
  \node[label, ubRed, anchor=north east] at (3.15,0.35)
    {badan ditempelkan\\pada sisi kiri bongkahan};
\end{scope}

% ---------- Panel B: sumbu tegak ----------
\begin{scope}[xshift=8.4cm]
  \node[judul, anchor=south] at (2.4,2.15)
    {2. Sumbu tegak diperiksa kemudian};

  \draw[bongkah] (0.6,0) rectangle (4.4,0.9);
  \node[label, anchor=north] at (2.5,-0.15) {bongkahan};

  % Posisi sebelum dan sesudah penyelesaian
  \draw[pemainlama] (1.75,0.9) rectangle (2.45,1.8);
  \draw[pemain] (1.75,0.45) rectangle (2.45,1.35);

  \draw[panah] (2.85,1.9) -- (2.85,1.3);
  \node[label, anchor=west] at (2.95,1.6) {jatuh};

  \draw[ubGreen, very thick] (0.6,0.9) -- (4.4,0.9);
  \node[label, ubGreen, anchor=south west] at (2.6,0.95)
    {kaki diletakkan\\pada permukaan};
\end{scope}
\end{tikzpicture}
}
\end{center}
\vspace{-0.04cm}
\begin{itemize}
  \item Bila seluruh perpindahan diperiksa sekaligus, program hanya tahu kedua
        persegi beririsan, tidak tahu dari sisi mana benda datang.
  \item Akibatnya muncul dua kesalahan khas: benda yang menabrak sisi
        bongkahan terangkat ke atasnya, atau benda yang jatuh di atasnya
        tertahan di udara.
  \item Penyelesaiannya: pindahkan pada sumbu mendatar saja lalu periksa dan
        perbaiki; setelah tuntas, baru pindahkan pada sumbu tegak.
  \item Sisi tabrakan dapat dipastikan dari arah gerak: bergerak ke kanan
        berarti menabrak dari kiri, dan sebaliknya.
\end{itemize}
\end{frame}

\begin{frame}[fragile]{Sumbu Mendatar Lebih Dahulu}
\lstinputlisting[language=Python,firstline=133,lastline=142]
  {../buku_ajar/code/python/ch07_tabrakan.py}
\kodejelas{Bila bergerak ke kanan, \texttt{rect.right} ditempelkan pada
\texttt{platform.rect.left}; bila bergerak ke kiri, \texttt{rect.left}
ditempelkan pada sisi kanannya. Perbaikan dilakukan pada \texttt{rect}, lalu
nilainya disalin kembali ke \texttt{position} agar posisi pecahan tetap menjadi
sumber kebenaran.}
\end{frame}

\begin{frame}[fragile]{Sumbu Tegak Sesudahnya}
\lstinputlisting[language=Python,firstline=144,lastline=156]
  {../buku_ajar/code/python/ch07_tabrakan.py}
\kodejelas{Saat jatuh, \texttt{rect.bottom} diletakkan pada permukaan bongkahan,
kecepatan tegak dinolkan, dan \texttt{on\_ground} diatur benar. Tanpa penolan
kecepatan, benda terus tertekan ke dalam bongkahan dan bergetar halus di
permukaannya. Gravitasi ditambahkan sebelum gerak tegak, mengikuti kaidah Euler
semi-implisit.}
\end{frame}

% ========================================================= bagian 4
\section{Kapan Persegi Pembatas Tidak Cukup}
\sectionframe{04}{Kapan Persegi Pembatas Tidak Cukup}
  {Ketelitian sampai titik gambar, dan harga yang harus dibayar}

\begin{frame}{Keterbatasan Persegi Pembatas}
\begin{center}
\adjustbox{max width=0.80\linewidth,max height=0.40\textheight}{%
  % chapters/07/figures/aabb-mask.tikz
% Kapan persegi pembatas tidak cukup: dua benda miring yang persegi
% pembatasnya beririsan, tetapi gambarnya tidak bersinggungan.
\begin{tikzpicture}[
  kotakpembatas/.style={draw=ubRed, thick, dashed},
  bentuk/.style={draw=ubPrimary, thick, fill=ubPrimary!20},
  panah/.style={-{Stealth[length=2.2mm]}, thick, draw=ubAccent},
  judul/.style={font=\footnotesize\bfseries, ubPrimary},
  label/.style={font=\scriptsize, align=center}
]

% ---------- Persegi pembatas beririsan ----------
\begin{scope}
  \node[judul, anchor=south] at (1.9,2.9) {diperiksa dengan persegi pembatas};

  \draw[bentuk] (0.35,0.35) -- (2.35,1.55) -- (0.35,1.55) -- cycle;
  \draw[bentuk] (1.45,0.35) -- (3.45,1.55) -- (1.45,1.55) -- cycle;
  \draw[kotakpembatas] (0.35,0.35) rectangle (2.35,1.55);
  \draw[kotakpembatas] (1.45,0.35) rectangle (3.45,1.55);

  \node[label, anchor=north, text width=52mm] at (1.9,-0.15)
    {Kedua persegi beririsan, sehingga program menyimpulkan kedua benda
     bersentuhan, padahal gambarnya tidak bersinggungan.};
\end{scope}

% ---------- Diperiksa dengan mask ----------
\begin{scope}[xshift=7.6cm]
  \node[judul, anchor=south] at (1.9,2.9) {diperiksa dengan \istilah{mask}};

  \draw[bentuk] (0.35,0.35) -- (2.35,1.55) -- (0.35,1.55) -- cycle;
  \draw[bentuk] (1.45,0.35) -- (3.45,1.55) -- (1.45,1.55) -- cycle;

  \node[label, anchor=north, text width=52mm] at (1.9,-0.15)
    {Yang dibandingkan adalah titik-titik yang benar-benar terisi. Karena tidak
     ada yang bertumpang tindih, kedua benda dinilai tidak bersentuhan.};
\end{scope}
\end{tikzpicture}
}
\end{center}
\vspace{-0.02cm}
\begin{itemize}
  \item Persegi pembatas selalu lebih besar daripada gambarnya, kecuali untuk
        benda yang memang berbentuk persegi.
  \item Pada benda miring atau berlubang, selisihnya dapat cukup besar sehingga
        dua benda dinilai bersentuhan padahal gambarnya masih berjarak.
  \item Gejalanya mudah dikenali: pemain merasa ``kena'' oleh sesuatu yang
        belum menyentuhnya, dan pada permainan yang menuntut ketepatan selisih
        beberapa piksel menentukan berhasil atau gagalnya percobaan.
\end{itemize}
\end{frame}

\begin{frame}{Pemeriksaan dengan \emph{Mask}}
\begin{itemize}
  \item Sebuah \emph{mask} adalah gambaran hitam-putih yang menyatakan titik
        mana saja pada permukaan yang benar-benar terisi.
  \item Mask dibuat dengan \texttt{pygame.mask.from\_surface(image)}, lalu
        dibandingkan melalui \texttt{overlap} dengan argumen kedua berupa
        \textbf{selisih letak} kedua benda, bukan letak salah satunya.
  \item Salah menentukan selisih letak itu merupakan sumber kekeliruan yang
        paling sering terjadi pada pemakaian mask.
  \item Buat mask sekali ketika benda dibuat, lalu simpan pada benda itu;
        membuatnya setiap bingkai membuang waktu karena menuntut pemeriksaan
        seluruh titik gambar.
\end{itemize}
\vspace{0.16cm}
\begin{kotak}
Pemeriksaan sampai titik gambar seperti ini disebut \emph{pixel-perfect}.
Namun pada program bab ini seluruh benda berbentuk persegi, sehingga persegi
pembatas sudah memadai dan mask hanya akan memperberat perhitungan.
\end{kotak}
\end{frame}

\begin{frame}{Memilih Metode yang Tepat}
{\scriptsize
\renewcommand{\arraystretch}{1.14}
\begin{tabularx}{\textwidth}{@{}>{\raggedright\arraybackslash}p{4.5cm}
    >{\raggedright\arraybackslash}p{3.3cm}
    >{\raggedright\arraybackslash}X@{}}
\toprule
\textbf{Keadaan} & \textbf{Metode yang dianjurkan} & \textbf{Alasan} \\
\midrule
Benda berbentuk persegi atau hampir persegi &
  Persegi pembatas (\emph{AABB}) &
  Selisih dengan bentuk sebenarnya kecil, sedangkan perhitungannya paling
  murah. \\
Banyak benda dan bergerak cepat &
  Persegi pembatas, dengan penyelesaian per sumbu &
  Ketelitian beberapa piksel tidak terasa, sedangkan mask memperberat setiap
  bingkai. \\
Benda miring, berlubang, atau bergerigi &
  \emph{Mask} &
  Persegi pembatasnya jauh lebih besar daripada gambar, sehingga tabrakan
  palsu sering terjadi. \\
Permainan yang menuntut ketepatan tinggi &
  \emph{Mask}, setelah pemeriksaan persegi pembatas &
  Persegi pembatas dipakai sebagai penyaring, lalu mask memeriksa hanya benda
  yang hampir bersentuhan. \\
\bottomrule
\end{tabularx}}
\vspace{-1.5cm}
\begin{kotak}
Cara yang paling sering dipakai pada permainan sungguhan adalah menggabungkan
keduanya: persegi pembatas diperiksa lebih dahulu karena murah, dan mask yang
mahal hanya dijalankan pada beberapa benda yang memang berdekatan.
\end{kotak}
\end{frame}

% ========================================================= bagian 5
\section{Program Tabrakan Lengkap}
\sectionframe{05}{Program Tabrakan Lengkap}
  {Tiga pemakaian tabrakan dalam satu prototipe yang dapat dimainkan}

\begin{frame}{Tiga Pemakaian Tabrakan di Dalamnya}
\awalan{Prototipe menyatukan gerak, jatuh, lompat, mendarat, terhalang,
mengumpulkan kepingan, dan kehilangan nyawa.}
\begin{itemize}
  \item \textbf{Bongkahan, dengan \texttt{colliderect} dan penyelesaian per
        sumbu.} Diperiksa di dalam \texttt{Player.update}; pemeriksaan ini
        bukan sekadar mengetahui, melainkan memperbaiki posisi dan kecepatan.
  \item \textbf{Kepingan, dengan \texttt{spritecollide} bernilai benar.}
        Kepingan yang tersentuh langsung dikeluarkan dari wadahnya, dan
        jumlahnya dihitung dari panjang daftar yang dikembalikan.
  \item \textbf{Rintangan, dengan \texttt{spritecollide} bernilai salah.}
        Rintangan tidak dikeluarkan karena harus tetap di tempatnya; sentuhan
        mengurangi satu nyawa dan mengembalikan pemain ke titik awal.
\end{itemize}
\vspace{0.16cm}
\begin{kotak}
Ketiganya memakai perintah dari bagian 2, tetapi dengan aturan permainan yang
berbeda. Inilah pekerjaan kedua yang tidak dapat dijawab oleh pustaka.
\end{kotak}
\end{frame}

\begin{frame}[fragile]{Pemeriksa Mandiri di Dalam Program}
\awalan{Penanganan tabrakan yang salah jarang menimbulkan pesan kesalahan.
Karena itu, kebenarannya diperiksa sendiri oleh program.}
\lstinputlisting[language=Python,firstline=186,lastline=193]
  {../buku_ajar/code/python/ch07_tabrakan.py}
\kodejelas{Fungsi menghitung berapa bongkahan yang badannya bertumpang tindih
dengan pemain. Nilainya seharusnya selalu nol: sesudah penyelesaian tabrakan,
pemain tidak boleh berada di dalam bongkahan mana pun. Nilai selain nol pada
suatu bingkai berarti penanganan tabrakan belum benar.}
\end{frame}

\begin{frame}[fragile]{Hasil Pemeriksaan Otomatis}
\begin{lstlisting}[language={},numbers=none]
Bingkai dijalankan       : 239
Kepingan terkumpul       : 1 dari 5
Skor                     : 10
Tersentuh rintangan      : 3
Bingkai bertumpang tindih: 0
\end{lstlisting}
\kodejelas{Baris terakhir adalah yang paling penting. Nilai nol menunjukkan
bahwa pada seluruh bingkai yang dijalankan, pemain tidak pernah berada di dalam
bongkahan---penanganan tabrakannya benar, walaupun pemain bergerak cepat,
melompat, dan terhalang berkali-kali.}
\vspace{0.12cm}
\begin{kotak}
Dengan pilihan \texttt{--auto-run}, program berjalan tanpa sentuhan manusia.
Ratusan bingkai dapat diperiksa dalam beberapa detik, pada beberapa laju
bingkai, tanpa harus memainkannya sendiri.
\end{kotak}
\end{frame}

% ========================================================= bagian 6
\section{Menganalisis Tabrakan yang Salah}
\sectionframe{06}{Menganalisis Tabrakan yang Salah}
  {Gejala, penyebab, perbaikan, dan kebiasaan yang mencegahnya}

\begin{frame}{Cara Menelusuri}
\awalan{Empat langkah berikut berlaku umum untuk kesalahan yang tidak
menimbulkan pesan kesalahan, bukan hanya untuk tabrakan.}
\begin{enumerate}
  \item \textbf{Dapatkan gejala yang dapat diulang.} Tentukan posisi awal, arah
        gerak, dan kecepatannya; gejala yang tidak dapat diulang tidak dapat
        diperbaiki.
  \item \textbf{Persempit wilayahnya.} Bila gejala hanya muncul ketika
        menyentuh bongkahan, pemeriksaan bongkahanlah yang diperiksa lebih
        dahulu.
  \item \textbf{Pasang alat ukur.} Tampilkan posisi sebelum dan sesudah
        pembaruan, nilai \texttt{rect}, dan kecepatan pada kedua sumbu.
  \item \textbf{Ubah satu hal pada satu waktu.} Mengubah dua hal sekaligus
        membuat Anda tidak tahu mana yang berpengaruh.
\end{enumerate}
\vspace{0.12cm}
\begin{kotak}
Pada program bab ini, penghitung tumpang tindih memainkan peran alat ukur itu.
\end{kotak}
\end{frame}

\begin{frame}{Tiga Gejala yang Paling Sering}
{\scriptsize
\renewcommand{\arraystretch}{1.10}
\begin{tabularx}{\textwidth}{@{}>{\raggedright\arraybackslash}p{3.9cm}
    >{\raggedright\arraybackslash}p{5.0cm}
    >{\raggedright\arraybackslash}X@{}}
\toprule
\textbf{Gejala} & \textbf{Penyebab} & \textbf{Perbaikan} \\
\midrule
Benda menembus bongkahan tipis, terutama saat bergerak cepat &
  Pada satu bingkai perpindahan melampaui tebal bongkahan, sehingga keduanya
  tidak pernah beririsan saat diperiksa (\emph{tunneling}). &
  Batasi panjang langkah, bagi perpindahan menjadi beberapa langkah kecil, atau
  tebalkan bongkahan. \\
Benda yang menabrak sisi bongkahan terangkat ke atasnya &
  Kedua sumbu diperiksa sekaligus, sehingga program tidak tahu sisi yang
  ditabrak dan selalu menyelesaikan pada sumbu tegak. &
  Pisahkan pemeriksaan menjadi dua tahap, mendahulukan sumbu mendatar. \\
Benda bergetar halus di permukaan bongkahan &
  Kecepatan tegak tidak dinolkan saat tabrakan diselesaikan, sehingga benda
  terus tertekan dan posisinya diperbaiki setiap bingkai. &
  Nolkan kecepatan pada sumbu yang diselesaikan. \\
Benda dinilai bersentuhan padahal gambarnya berjarak &
  Benda miring atau berlubang, sehingga persegi pembatasnya jauh lebih besar
  daripada gambarnya. &
  Pakai \emph{mask}, dengan persegi pembatas sebagai penyaring lebih dahulu. \\
\bottomrule
\end{tabularx}}
\vspace{0.12cm}
\begin{kotak}
Ketiga gejala itu akibat pilihan cara, bukan kesalahan penulisan, sehingga
tidak dapat ditemukan oleh pemeriksa sintaksis. Kemampuan menganalisis seperti
ini membedakan program yang ``berjalan'' dari program yang ``benar''.
\end{kotak}
\end{frame}

\begin{frame}{Kebiasaan yang Mencegahnya}
\begin{itemize}
  \item \textbf{Selesaikan satu sumbu sampai tuntas} sebelum menyentuh sumbu
        berikutnya. Ini kebiasaan yang paling banyak mencegah kesalahan.
  \item \textbf{Tampilkan keadaan yang Anda curigai} pada layar atau terminal;
        angka yang terlihat menjawab pertanyaan yang tidak dapat dijawab oleh
        gambar.
  \item \textbf{Sediakan pemeriksa mandiri di dalam program}, yang menjawab
        ``apakah keadaan terlarang sedang terjadi?'' pada setiap bingkai.
  \item \textbf{Jalankan program pada beberapa kecepatan}; kesalahan
        \emph{tunneling} dan pembulatan hanya muncul pada kecepatan tertentu.
\end{itemize}
\vspace{0.14cm}
\begin{kotak}
Bab ini menutup CPMK 1: perkakas Blok I---lingkungan kerja, siklus permainan,
gerak, fisika, dan tabrakan---telah lengkap.
\end{kotak}
\end{frame}

% --------------------------------------------------------- aktivitas
\begin{frame}{Aktivitas di Kelas \textperiodcentered\ Praktikum 05}
\awalan{Pengalaman belajar minggu 5 pada RPS: mengumpulkan kepingan dan
menghindari rintangan, dengan skor yang bertambah setiap kali kepingan
tersentuh.}
\begin{enumerate}
  \item Ketik ulang \texttt{ch07\_tabrakan.py} ke folder \texttt{code}, lalu
        mainkan sampai seluruh kepingan terkumpul.
  \item Jalankan dengan \texttt{--frames 300 --auto-run --fps 30}, \texttt{60},
        dan \texttt{120}; catat jumlah bingkai bertumpang tindih pada tiap laju.
  \item Rusak dengan sengaja: satukan kedua sumbu pada \texttt{update} pemain,
        lalu amati apa yang terjadi saat menabrak sisi bongkahan.
  \item Kembalikan program, lalu hapus penolan kecepatan tegak sesudah
        mendarat; perhatikan perilaku pemain di atas bongkahan.
  \item Perbaiki kembali, ubah \texttt{MOVE\_SPEED} menjadi 1500, dan periksa
        apakah penghitung tumpang tindih masih nol.
  \item Tambahkan satu bongkahan dan dua kepingan, tulis catatan ketiga gejala
        beserta perbaikannya, lalu simpan ke repositori Git.
\end{enumerate}
\vspace{0.10cm}
{\footnotesize\color{ubAccent}\itshape Rincian lengkap: Praktikum 05 pada Bab VII
buku ajar. Kumpulkan melalui repositori Git pribadi.}
\end{frame}

% --------------------------------------------------------- rangkuman
\begin{frame}{Rangkuman}
\begin{itemize}
  \item Kelas turunan \texttt{pygame.sprite.Sprite} menyatukan gambar, persegi
        pembatas, dan perilaku benda; \texttt{image} dan \texttt{rect} wajib
        bernama demikian.
  \item \texttt{pygame.sprite.Group} mengerjakan \texttt{update} dan
        \texttt{draw} untuk seluruh anggotanya; satu benda boleh berada pada
        beberapa wadah.
  \item Tiga perintah: \texttt{colliderect} untuk dua persegi,
        \texttt{spritecollide} untuk satu benda terhadap wadah, dan
        \texttt{groupcollide} untuk dua wadah.
  \item Mendeteksi dan menyelesaikan adalah dua pekerjaan berbeda; tabrakan
        diselesaikan sumbu demi sumbu, mendahulukan sumbu mendatar.
  \item Persegi pembatas kurang teliti untuk benda miring atau berlubang;
        \emph{mask} dipakai setelah persegi pembatas menyaring kandidat.
  \item Kesalahan tabrakan tidak menimbulkan pesan kesalahan; pemeriksa mandiri
        yang selalu bernilai nol menemukannya lebih cepat daripada mata.
\end{itemize}
\end{frame}

% --------------------------------------------------------- cek
\begin{frame}{Cek Pemahaman}
\begin{enumerate}
  \item Mengapa dua persegi yang hanya bersentuhan pada tepinya dinilai tidak
        beririsan, dan mengapa sifat itu berguna bagi benda yang telah
        ditempelkan pada bongkahan?
  \item Mengapa sumbu mendatar harus diselesaikan lebih dahulu, dan gejala apa
       yang muncul bila kedua sumbu diperiksa sekaligus?
  \item Kapan persegi pembatas tidak memadai sehingga \emph{mask} diperlukan,
        dan mengapa persegi pembatas tetap dipakai lebih dahulu?
\end{enumerate}
\vspace{0.35cm}
{\footnotesize\color{ubMuted}\itshape Jawab tanpa membuka catatan, lalu
bandingkan dengan pasangan di sebelah Anda.}
\end{frame}

% --------------------------------------------------------- penutup
\penutup{Tugas Praktikum 05\\ Kuis 1 (Materi 1--4)}
  {Materi 6 membuka Blok II: NPC yang mengambil keputusan sendiri. Anda akan
   menghitung vektor arah menuju pemain, menormalisasinya, membatasi besar
   kecepatan, dan membandingkan AI reaktif dengan AI berbasis perencanaan.}

\end{document}
